\documentclass[11pt]{article}

\usepackage[margin=1in]{geometry}
\usepackage{amsmath,amssymb}
\usepackage{booktabs}
\usepackage{graphicx}
\usepackage{xcolor}
\usepackage{listings}
\usepackage{lmodern}
\usepackage[T1]{fontenc}
\usepackage{microtype}
\usepackage[hidelinks]{hyperref}
\usepackage[capitalise]{cleveref}
\usepackage{natbib}

\lstset{
  basicstyle=\ttfamily\small,
  breaklines=true,
  columns=fullflexible,
  frame=single,
  framesep=4pt,
  xleftmargin=6pt,
  literate={~}{{\textasciitilde}}1
}

\newcommand{\pass}{\textsc{pass}}
\newcommand{\usable}{\textsc{usable}}
\newcommand{\vuln}{\textsc{vulnerable}}
\newcommand{\dfaeq}{\textsc{dfa-eq}}
\newcommand{\exact}{\textsc{exact}}

\title{\textbf{Passing Is Not Shipping}\\[0.4em]
\large Correctness, Semantic Equivalence and ReDoS Safety in\\
LLM-Generated Regular Expressions, and the Limits of Measuring Them}

\author{Foothills Labs\\\texttt{info@foothills-labs.com}}

\date{August 2026}

\begin{document}
\maketitle

\begin{abstract}
\noindent
We evaluate eleven contemporary language models on natural-language-to-regex
generation over 450 tasks with three samples each ($14{,}850$ requests,
\$10.95), scoring every answer on three axes: whether it satisfies the task's
tests, whether it denotes the same language as a human reference, and whether
it is vulnerable to regular-expression denial of service (ReDoS).

Our headline result is a negative one about a composite metric we
constructed and then had to abandon. A composite requiring correctness,
safety and semantic equivalence appears to show that ``passing is twice as
easy as shipping'': models pass 38.0--47.5\% of tasks but satisfy all three
criteria on only 17.1--24.8\%. Decomposing that gap shows that
\textbf{87\% of it is contributed by the semantic-equivalence term}, and a
manual audit shows that term is itself dominated by reference-set error and
task underspecification rather than model behaviour. Using only the two
criteria that never consult a human answer key, the gap is
\textbf{2.9 percentage points}: just 7.4\% of functionally correct patterns
are ReDoS-vulnerable.

That figure is itself informative by contrast. Joint correctness-and-security
benchmarks in general code generation report far larger penalties. BaxBench
finds roughly half of correct backends exploitable. SecureAgentBench reports
15.2\% correct-and-secure for its best agent. The regex domain does \emph{not}
reproduce that pattern, and we report the difference rather than the
similarity. Separately, \emph{models are safer than the humans whose answers
define the benchmark}: 13.6\% of the corpus's human-written reference patterns
are ReDoS-vulnerable against 9.0\% for models pooled, and every individual
model is safer than the reference set.

The audit behind this is the paper's methodological contribution. Metrics
that compare generated output against a human answer key inherit that key's
error rate. In a random sample of disagreements where a model passed every
test yet was scored as semantically different, the model was clearly at fault
in roughly 15\% of cases. The rest split between demonstrably incorrect
reference patterns (including a reference character class containing
a literal \texttt{|}) and prompts that never specified the property in
dispute. Because that unreliable term supplies 87\% of our composite's
headline gap, the composite was measuring its own weakest component. We also show that independent binomial intervals, the default in this
literature, are the wrong inferential tool when all systems answer identical
items, and that a paired bootstrap resolves nine times as
many pairwise comparisons on the same data.

We therefore decline to publish a ranking. We release all raw generations,
scoring code and adjudications, and recommend that evaluations separate
reference-dependent from reference-independent metrics and assign them
different evidentiary weight.
\end{abstract}

\vspace{0.5em}
\noindent\textbf{Artifacts.} All model outputs, scoring code, per-case
adjudications and analysis scripts are available at
\url{https://github.com/foothills-labs/regexleaderboard}. Every number in
this paper recomputes offline from committed data with no API access.

\section{Introduction}

Regular expressions make an unusually good probe for language-model
evaluation. The task is compact. The output runs. Correctness is partly
decidable, which is rare. And a syntactically valid, test-passing answer can
still be a security defect, which is what drew us to the task in the first
place.

The standard protocol for this task scores a candidate against a reference
by DFA equivalence~\citep{locascio2016neural,kushman2013semantic}, sometimes
alongside test-based correctness. Neither criterion notices catastrophic
backtracking. A pattern can satisfy every provided example, denote exactly the
reference language, and still admit inputs on which matching takes time
exponential in the input. That is the ReDoS vulnerability
class~\citep{davis2018redos}, a denial-of-service vector in deployed systems
and common enough in human-written code to have prompted ecosystem-scale
study.

\subsection{Contributions}

\begin{enumerate}
  \item \textbf{A three-axis evaluation of eleven current models}
  (\cref{sec:results}), reporting test correctness, semantic equivalence and
  ReDoS safety jointly, together with a composite (\usable{}) requiring all
  three. Joint correctness-and-security scoring is established prior art
  (\cref{sec:jointprior}). What we add is its transfer to
  natural-language-to-regex generation, a setting that differs from existing
  joint benchmarks in three respects: the vulnerability class is ReDoS rather
  than CWE-classified defects, correctness is \emph{partly decidable} rather
  than test-only, and a per-task human reference exists, which enables
  \cref{sec:humanbaseline}.

  \item \textbf{A human baseline for ReDoS} (\cref{sec:humanbaseline}). We
  apply the identical safety screen to the corpus's own reference patterns.
  Models are uniformly safer than the reference set, which relocates the
  finding from a claim about model deficiency to a claim about the
  distribution models were trained on.

  \item \textbf{A quantified audit of reference-standard error}
  (\cref{sec:goldstandard}). We sample disagreements and adjudicate them,
  finding that the majority of apparent semantic errors are not model errors.
  Existing joint benchmarks are largely reference-free. They execute tests and
  exploits instead of comparing against a gold artifact, so the failure mode
  does not arise for them. It arises here because the NL-to-regex literature is
  built on reference comparison.

  \item \textbf{A methodological correction} (\cref{sec:paired}). We show
  that independent intervals substantially understate resolvable differences
  when all systems answer identical items, and give the paired alternative.

  \item \textbf{Full artifact release} including raw generations, which
  permits any of our scoring decisions to be overturned without re-running
  inference.
\end{enumerate}

\subsection{What this paper does not claim}

We do not publish a ranking of the eleven models. \cref{sec:power} shows
that 62\% of tasks yield identical outcomes across all eleven systems, and
that only nine of fifty-five pairwise comparisons resolve at 95\% confidence
even under the correct paired test. Reporting an ordinal leaderboard from
this data would be unsupported by it.

We also do not claim novelty for DFA-equivalence checking, which is standard
in this literature~\citep{locascio2016neural}, nor for the scoring engine,
which is prior work by the present authors released
separately~\citep{regexbench} and used here as a pinned dependency.

\section{Background and Related Work}

\paragraph{NL-to-regex generation.}
The task was established by \citet{kushman2013semantic} (KB13) and extended
by \citet{locascio2016neural} (NL-RX), both of which adopt DFA equivalence
against a reference as the correctness criterion. That is the right choice in principle. $\texttt{[0-9]+}$ and
$\texttt{[0-9][0-9]*}$ denote the same language and differ as strings, so
text comparison would mark one of them wrong. But the criterion inherits any
defect in the reference set, which is the dependency this paper quantifies.

\paragraph{ReDoS.}
Catastrophic backtracking arises when a backtracking matcher can decompose a
prefix in exponentially many ways, canonically via a quantifier applied to a
quantified group, as in \texttt{(a+)+}. \citet{davis2018redos} document its prevalence at ecosystem scale, and
\citet{davis2019linguafranca} study regex portability and re-use across
languages. \citet{icpc2024llmredos} examine
ReDoS specifically in LLM-generated patterns and report a skew toward polynomial over exponential vulnerability. We compare
against that finding in \cref{sec:humanbaseline}.

\subsection{Joint correctness-and-security benchmarks}
\label{sec:jointprior}

Evaluating functional correctness and security together is an active area, and
this paper is not its first instance. \citet{peng2025cweval} introduce CWEval,
scoring 119 security-critical tasks across 31 CWEs and five languages with
outcome-driven oracles for both properties, and report a substantial
population of functionally correct but insecure generations.
\citet{vero2025baxbench} present BaxBench, 392 backend tasks with expert
exploits, and find that roughly half of functionally correct solutions are
exploitable. \citet{chen2025secureagentbench} evaluate coding \emph{agents} on
105 repository-level tasks, reporting 15.2\% correct-and-secure for the best
agent and 9.2\% on average. \citet{patir2025dualgauge} automate the joint
pipeline in DualGauge over 154 tasks and 10 models, observing secure-pass@1
below 12\% while functional pass@1 exceeds 50\%.

What we do here differs in domain, not in framing. The artifact is a single
expression instead of a program. The vulnerability class is ReDoS instead of a
CWE taxonomy. Correctness is partly decidable, not established by tests alone.
And the NL-to-regex literature supplies a human reference for every task,
which gives us a human baseline on the same instrument. That last point is
what makes \cref{sec:humanbaseline} possible. Reference-free benchmarks cannot
run that comparison at all.

Our numbers turned out not to match theirs. \cref{sec:crossbench} works
through the difference.

\paragraph{Sampling-based metrics.}
We use the unbiased $\pass{}@k$ estimator of \citet{chen2021codex}, with the
degeneracy noted in \cref{sec:metrics}.

\paragraph{Serving-layer variance.}
Hosted inference routers may serve the same model slug from multiple
providers at differing numerical precision. \citet{openrouterpinning} report
that 31 of 32 surveyed AI-safety codebases using such a router did not pin
the serving provider, making their results a measurement of the router's
routing policy in addition to the model. \cref{sec:serving} describes the
discipline we adopt in response.

\section{Method}

\subsection{Task and prompt}

Each task supplies a natural-language description. The model returns one
regular expression. The prompt is fixed across all models and all tasks:

\begin{quote}\small
\textbf{system:} You translate a natural-language description into a single
Python \texttt{re}-compatible regular expression. Reply with ONLY the pattern
in one fenced code block, no explanation.\\[0.3em]
\textbf{user:} \textit{\{task description, verbatim\}}\\[0.3em]
Reply with only the regular expression pattern in a single fenced code block.
\end{quote}

We do not tune this prompt. Prompt optimisation would raise all scores and
destroy comparability with prior work on the same corpus. \cref{sec:threats}
notes the corresponding threat.

\paragraph{Extraction.} We take the final fenced code block, else the whole
trimmed reply, then strip at most one layer of host-language quoting
(\texttt{r'\ldots'}, \texttt{'\ldots'}, \texttt{"\ldots"}, backticks,
\texttt{/\ldots/flags}). Scored literally, \texttt{r'\textbackslash d+\$'} is
a pattern matching the letter \texttt{r}, and would fail for reasons
unrelated to regex competence. This affected 7--18 of 1{,}350 responses per
model, under 1.4\% throughout, and changes no conclusion. We release both the
normalised and the unnormalised scores.

\subsection{Corpus}

We use Re(gEx$|$DoS)Eval~\citep{regexeval}, 762 tasks drawn from real user
requests, each with a description, positive and negative test strings, and a
human-written reference pattern. We evaluate 450 tasks, selected at uniform
stride across the corpus rather than as a prefix (the corpus is
difficulty-ordered at its head). The subset size was budget-determined.

Match semantics are set per corpus. Re(gEx$|$DoS)Eval references satisfy
their own tests at 100\% under search semantics and 94\% under full-match.
Choose wrong and 46 gold patterns are scored as failures. We verified
the loaded configuration by scoring every reference against its own tests
before the run (\cref{sec:goldstandard}).

\subsection{Models and serving discipline}
\label{sec:serving}

Eleven models spanning frontier and open-weights systems across a $98\times$
price range. Every request pins a single named provider with fallback
disabled, so a request is either served by the named endpoint or fails
visibly. Substitution failures did occur and are reported as failures rather
than silently re-routed (\cref{app:failures}).

For open-weights models the pin includes quantisation where the router
discloses it: GLM-5.2 and DeepSeek-V4-Flash are served at 4-bit by some
providers and 8-bit by others, and we pin 8-bit. One pin went stale within an
hour of being set (the provider ceased serving that model), so pins are
re-verified immediately before each run.

The pin is a request, not a guarantee. So we record which provider actually
served each response and check afterwards that it matches. All eleven models
came back served entirely by the endpoint they were pinned to.

\subsection{Metrics}
\label{sec:metrics}

Let $T=(d, P, N, r)$ be a task with description $d$, positive examples $P$,
negative examples $N$ and reference pattern $r$. Let $p$ be a candidate.

\paragraph{Correctness.}
\[
\mathrm{correct}(p,T) \;=\; \big[\forall x \in P:\ \mathrm{match}(p,x)\big]
\;\wedge\;
\big[\forall y \in N:\ \neg\,\mathrm{match}(p,y)\big]
\]
evaluated with CPython's \texttt{re} under the corpus's declared semantics.
This is \emph{reference-independent}.

\paragraph{Semantic equivalence.}
Both $p$ and $r$ are compiled to finite automata and compared, yielding a
verdict in $\{\textsc{eq}, \textsc{diff}, \textsc{undec}, \textsc{unsup}\}$.
When the verdict is \textsc{diff} the engine returns a \emph{witness}: a
shortest string in the symmetric difference, which renders the judgment
falsifiable by inspection.

Equivalence is decidable only on the regular subset. Backreferences make the
language non-regular, and equivalence then has no algorithmic answer at all.
Those cases return \textsc{undec}. Between 82 and 133 of 450 tasks per model
land there.
We therefore report two figures:
\[
\dfaeq{} = \Pr[\textsc{eq}], \qquad
\dfaeq{}_{\text{dec}} = \Pr[\textsc{eq} \mid \text{verdict} \neq \textsc{undec}].
\]
The first counts undecidable comparisons as failures. It is a lower bound
and cannot flatter. The second isolates the model from the engine's reach.
Report only the second and you inflate. Report only the first and you charge
the model for a theorem. This metric is \emph{reference-dependent}, which
\cref{sec:goldstandard} shows to be decisive.

\paragraph{ReDoS safety.}
$\mathrm{vuln}(p) = [\mathrm{screen}(p) \neq \textsc{safe}]$, where
$\mathrm{screen}$ performs structural analysis for known-vulnerable shapes
followed by an empirical pass against attack strings under timeout. The
structural pass covers three of the five families catalogued
by~\citet{icpc2024llmredos}. So \vuln{} is a \emph{lower bound}, and
\textsc{safe} denotes a screening outcome, not a proof. This metric is
\emph{reference-independent}.

\paragraph{Composite.}
\[
\mathrm{usable}(p,T) = \mathrm{correct}(p,T)
  \;\wedge\; \neg\,\mathrm{vuln}(p)
  \;\wedge\; \big[\text{verdict}(p,r) \neq \textsc{diff}\big]
\]
The third conjunct excludes only \emph{proven} difference. Undecidable and
unsupported verdicts do not disqualify a pattern. \usable{} is
reference-dependent through that conjunct, and inherits the error
characterised in \cref{sec:goldstandard}.

\paragraph{Sampling.}
With $n$ samples per task of which $c$ satisfy a predicate, we use the
unbiased estimator of~\citet{chen2021codex}:
\[
\widehat{\mathrm{metric}}@k \;=\;
\mathbb{E}_{T}\left[\,1 - \frac{\binom{n-c}{k}}{\binom{n}{k}}\right].
\]
We draw $n=3$ and report $k=3$. Note that at $n=k$ this estimator degenerates to $\mathbf{1}[c \geq 1]$ and
buys no variance reduction over the naive any-of-$n$ statistic. We report it as
$@3$ for comparability, not because the estimator is doing any work at this
setting.
$n>k$ would be required for the estimator's variance properties to apply,
and is left to future work with a larger budget.

\subsection{Sampling configuration}

\paragraph{Reasoning disabled.} Four of the eleven models emit no hidden
reasoning tokens. The other seven do. Evaluating them together with reasoning enabled
would confound model identity with inference-time compute budget. All
requests therefore disable reasoning, and reasoning-token counts are recorded
per response to verify compliance (observed: zero throughout).

\paragraph{Temperature unset.} We set
\texttt{require\_parameters=true}, which restricts routing to providers that
honour submitted sampling parameters rather than silently discarding them.
Six of the eleven models reject \texttt{temperature} outright. Together, the
two settings make the request unroutable. We therefore omit \texttt{temperature} entirely and
each model samples at its provider default. That keeps the non-silent-discard guarantee and costs us controlled
sampling. We think it is the right trade and note it as a limitation.

Since $k>1$ is only informative under output diversity, we verified diversity
directly: repeated queries on a fixed task yielded 3, 3, 3 and 2 distinct
patterns for \texttt{claude-opus-5}, \texttt{gpt-5.6-sol}, \texttt{kimi-k3}
and \texttt{glm-5.2} respectively. The analysis pipeline additionally warns
if $>90\%$ of a model's tasks return identical samples.

\subsection{Controls}

Three synthetic responses go through the identical scoring path on every run.
The task's own reference must come back correct and usable. The pattern
\texttt{z\{5\}} must fail every correctness check. \texttt{(a+)+b} must be
flagged vulnerable. If any control misbehaves we throw the run away instead of
reporting it.

The failure this catches is a scorer that quietly returns nulls, which looks
exactly like every model performing badly. All controls behaved on all eleven
models.

\section{Results}
\label{sec:results}

\subsection{Main results}

\begin{table}[htbp]
\centering
\footnotesize
\setlength{\tabcolsep}{3pt}
\begin{tabular}{lrrrrrrrr}
\toprule
Model & \pass{}@3 & \usable{}@3 & \vuln{}@3 & \dfaeq{}@3 & \dfaeq{}$_{\text{dec}}$ & \exact{}@3 & undec. & fail \\
\midrule
\texttt{kimi-k3} & 47.2 & 24.8 & 14.1 & 15.2 & 17.5 & 6.3 & 82 & 16 \\
\texttt{claude-opus-5} & 47.5 & 23.0 & 15.5 & 14.0 & 17.0 & 5.4 & 109 & 36 \\
\texttt{qwen3.6-max-preview} & 42.4 & 21.6 & 9.1 & 13.8 & 17.4 & 3.8 & 94 & 0 \\
\texttt{gpt-5.6-sol} & 42.2 & 21.1 & 10.2 & 9.6 & 13.6 & 2.0 & 133 & 1 \\
\texttt{deepseek-v4-flash-0731} & 38.0 & 19.8 & 12.0 & 11.3 & 14.3 & 3.3 & 93 & 0 \\
\texttt{qwen3.6-plus} & 39.8 & 19.8 & 9.8 & 11.6 & 14.8 & 3.8 & 99 & 0 \\
\texttt{glm-5.2} & 42.4 & 18.7 & 14.2 & 10.4 & 12.9 & 3.8 & 86 & 0 \\
\texttt{gpt-5.6-luna} & 39.3 & 18.7 & 11.8 & 9.3 & 12.2 & 2.0 & 113 & 1 \\
\texttt{gpt-5.6-terra} & 42.2 & 18.7 & 12.0 & 10.2 & 13.1 & 2.0 & 100 & 0 \\
\texttt{claude-sonnet-5} & 40.7 & 18.0 & 10.9 & 10.4 & 13.3 & 3.3 & 97 & 0 \\
\texttt{gemini-3.1-flash-lite} & 38.7 & 17.1 & 12.0 & 9.3 & 11.8 & 3.1 & 95 & 0 \\
\bottomrule
\end{tabular}

\caption{Primary results, 450 tasks, $n=k=3$. All figures percentages except
\emph{undec.} (count of tasks with undecidable equivalence, of 450) and
\emph{fail} (count of failed requests, of 1{,}350). Ordering is by
\usable{}@3 and should be read as a banding, not a ranking
(\cref{sec:power}).}
\label{tab:main}
\end{table}

The apparent headline in \cref{tab:main} is the gap between $\pass{}@3$
(38.0--47.5\%) and $\usable{}@3$ (17.1--24.8\%): approximately half of all
test-passing generations fail at least one of the remaining two criteria, and
the ratio is stable across the entire price range.

\cref{sec:decomp} shows that this reading is wrong, or at least that it
attributes the gap to the wrong cause. We present it here as stated because it
is what the composite reports, and because constructing such a composite is a
natural thing to do. The point of the next section is that building one
without decomposing it was a mistake, and it was ours.

\subsection{Decomposing the gap}
\label{sec:decomp}

\usable{} conjoins three conditions. \cref{tab:decomp} attributes the drop
from \pass{}@3 to \usable{}@3 between the safety term and the equivalence
term, and additionally reports the rate of ReDoS vulnerability
\emph{conditional on functional correctness}.

\begin{table}[htbp]
\centering
\footnotesize
\setlength{\tabcolsep}{4pt}
\begin{tabular}{lrrrrrr}
\toprule
Model & \pass{}@3 & $-$safety & C\&S@3 & $-$equiv. & \usable{}@3 & vuln.$\mid$correct \\
\midrule
\texttt{kimi-k3} & 46.5 & 2.5 & 44.0 & 20.2 & 23.8 & 5.6 \\
\texttt{claude-opus-5} & 46.1 & 4.9 & 41.2 & 20.4 & 20.8 & 10.0 \\
\texttt{qwen3.6-max-preview} & 42.4 & 2.7 & 39.8 & 18.2 & 21.6 & 7.0 \\
\texttt{gpt-5.6-sol} & 42.1 & 1.6 & 40.5 & 19.6 & 20.9 & 5.6 \\
\texttt{deepseek-v4-flash-0731} & 38.0 & 1.8 & 36.2 & 16.4 & 19.8 & 5.7 \\
\texttt{qwen3.6-plus} & 39.8 & 3.3 & 36.4 & 16.7 & 19.8 & 8.1 \\
\texttt{glm-5.2} & 42.4 & 4.0 & 38.4 & 19.8 & 18.7 & 9.3 \\
\texttt{gpt-5.6-luna} & 39.2 & 2.9 & 36.3 & 17.8 & 18.5 & 8.1 \\
\texttt{gpt-5.6-terra} & 42.2 & 2.9 & 39.3 & 20.7 & 18.7 & 7.3 \\
\texttt{claude-sonnet-5} & 40.7 & 2.7 & 38.0 & 20.0 & 18.0 & 6.0 \\
\texttt{gemini-3.1-flash-lite} & 38.7 & 2.9 & 35.8 & 18.7 & 17.1 & 8.3 \\
\midrule
\textit{mean} & --- & 2.9 & --- & 18.9 & --- & 7.4 \\
\bottomrule
\end{tabular}

\caption{Decomposition of the \pass{}$\to$\usable{} gap. ``C\&S'' is
correct-and-secure: correctness conjoined with safety, omitting the
equivalence term. The final column is the share of functionally correct
patterns that are ReDoS-vulnerable.}
\label{tab:decomp}
\end{table}

Safety costs 2.9 percentage points on average. Equivalence costs 18.9. So the
equivalence term supplies \textbf{87\%} of the gap that made the composite
look interesting.

That would be unremarkable if the equivalence term were sound. It is not.
\cref{sec:goldstandard} shows, from an independent audit, that roughly 85\% of
what it counts against a model is reference-set error or an underspecified
prompt. Our composite was mostly reporting its own weakest component back to
us.

We did not see this coming. The decomposition was run late, after the audit
and after the first draft, because a reviewer of an earlier version asked how
our numbers compared to joint benchmarks in other domains. Answering that
required computing correct-and-secure without the equivalence term, and the
answer changed the paper.

Restricting attention to the two reference-independent criteria, the finding
is far more modest and far better supported: \textbf{7.4\% of functionally
correct regular expressions are ReDoS-vulnerable} (range 5.6--10.0\%).

We count 135 model-task pairs in which a pattern satisfied every test and was
simultaneously screened as unsafe. As a rate over correct generations instead of a raw count, that is the 7.4\%
above. About one in fourteen, not one in two.

$\exact{}@3$, string identity with the reference, runs from 2.0\% to 6.3\%.
String-comparison scoring would therefore misrank essentially every system,
which is what justifies the automata-theoretic treatment.

\subsection{Sample budget and the fragility of ordering}
\label{sec:atk}

Because $n=k=3$ renders the $\pass{}@k$ estimator degenerate
(\cref{sec:metrics}), we additionally report $@1$, which is the per-sample
success rate and the quantity most directly comparable to single-sample
protocols. \cref{tab:at1} gives both.

\begin{table}[htbp]
\centering
\small
\begin{tabular}{lrrrrrr}
\toprule
Model & \pass{}@1 & \pass{}@3 & \usable{}@1 & \usable{}@3 & \vuln{}@1 & $n{<}3$ \\
\midrule
\texttt{kimi-k3} & 35.6 & 46.5 & 16.0 & 23.8 & 9.2 & 6 \\
\texttt{claude-opus-5} & 41.4 & 46.1 & 17.3 & 20.8 & 10.7 & 12 \\
\texttt{qwen3.6-max-preview} & 37.0 & 42.4 & 16.5 & 21.6 & 7.4 & 0 \\
\texttt{gpt-5.6-sol} & 37.3 & 42.1 & 16.0 & 20.9 & 8.4 & 1 \\
\texttt{deepseek-v4-flash-0731} & 30.0 & 38.0 & 14.4 & 19.8 & 7.7 & 0 \\
\texttt{qwen3.6-plus} & 35.5 & 39.8 & 15.9 & 19.8 & 8.0 & 0 \\
\texttt{glm-5.2} & 34.1 & 42.4 & 13.5 & 18.7 & 9.6 & 0 \\
\texttt{gpt-5.6-luna} & 33.8 & 39.2 & 15.2 & 18.5 & 8.7 & 1 \\
\texttt{gpt-5.6-terra} & 36.6 & 42.2 & 14.8 & 18.7 & 8.9 & 0 \\
\texttt{claude-sonnet-5} & 36.8 & 40.7 & 15.3 & 18.0 & 9.3 & 0 \\
\texttt{gemini-3.1-flash-lite} & 33.2 & 38.7 & 13.9 & 17.1 & 9.5 & 0 \\
\bottomrule
\end{tabular}

\caption{Per-sample ($@1$) against any-of-three ($@3$) estimates, computed
from per-sample success counts. The final column gives the number of tasks
with fewer than three usable samples, which are excluded from that model's
$@3$ estimate. That is why $@3$ here differs slightly from \cref{tab:main},
which is computed over all answered tasks.}
\label{tab:at1}
\end{table}

The comparison is instructive beyond bookkeeping. At $@3$,
\texttt{kimi-k3} attains the highest \usable{} at 23.8\%. At $@1$ the leader
is \texttt{claude-opus-5} on 17.3\%, and \texttt{kimi-k3} drops to joint
third.
\emph{The induced ordering depends on the sample budget.} Systems differ not
only in per-attempt quality but in how much they benefit from additional
attempts, and a leaderboard that fixes $k$ silently fixes a weighting between
those two properties. This is a further reason to prefer banded reporting
(\cref{sec:power}), and a reason to state $k$ prominently whenever an @$k$
metric is ranked.

\subsection{ReDoS and the human baseline}
\label{sec:humanbaseline}

\cref{tab:vuln} applies the identical safety screen to the corpus's 450
human-written reference patterns. For comparability with the human set (one
pattern per task) we score one sample per task per model rather than the
any-of-three statistic of \cref{tab:main}.

\begin{table}[htbp]
\centering
\small
\begin{tabular}{lrrrr}
\toprule
Source & $n$ & vulnerable (\%) & exponential (\%) & polynomial (\%) \\
\midrule
\textit{Human reference answers} & 450 & 13.6 & 6.4 & 7.1 \\
\midrule
\texttt{kimi-k3} & 447 & 10.7 & 5.4 & 5.4 \\
\texttt{claude-opus-5} & 444 & 10.6 & 7.4 & 3.2 \\
\texttt{glm-5.2} & 450 & 10.2 & 5.8 & 4.4 \\
\texttt{gemini-3.1-flash-lite} & 450 & 9.8 & 5.1 & 4.7 \\
\texttt{gpt-5.6-sol} & 450 & 9.3 & 5.1 & 4.2 \\
\texttt{claude-sonnet-5} & 450 & 8.9 & 6.0 & 2.9 \\
\texttt{gpt-5.6-terra} & 450 & 8.7 & 4.9 & 3.8 \\
\texttt{qwen3.6-plus} & 450 & 8.4 & 5.1 & 3.3 \\
\texttt{gpt-5.6-luna} & 450 & 8.0 & 4.0 & 4.0 \\
\texttt{deepseek-v4-flash-0731} & 450 & 7.3 & 4.4 & 2.9 \\
\texttt{qwen3.6-max-preview} & 450 & 7.3 & 4.7 & 2.7 \\
\midrule
\textit{All models pooled} & 4941 & 9.0 & 5.3 & 3.8 \\
\bottomrule
\end{tabular}

\caption{ReDoS screening of model outputs against the corpus's own
human-written reference patterns, one pattern per task in all cases.}
\label{tab:vuln}
\end{table}

Every model is safer than the reference set. This inverts the expected
framing. Generation does not introduce ReDoS vulnerability relative to human
practice. ReDoS vulnerability is endemic to human-written regular expressions,
at a rate above what these models reproduce.

We regard this as the paper's most consequential empirical finding, because
it changes the implied remedy. A model-deficiency framing implies waiting for
better models. A distributional framing implies that safety screening must be
applied at the point of use, since neither the generator nor the human corpus
it was trained on performs that check.

\paragraph{Relation to prior work.} \citet{icpc2024llmredos} report that
LLM-generated regexes skew toward polynomial ReDoS. We do not reproduce this
cleanly: pooled across models we observe 5.3\% exponential against 3.8\%
polynomial, the opposite direction, while the human references skew
polynomial (6.4\% vs.\ 7.1\%). The subcategory counts are small, so we make no strong claim. The
discrepancy is worth a replication.

\subsection{Comparison against joint benchmarks in other domains}
\label{sec:crossbench}

Because \usable{} carries a semantic-equivalence conjunct that the joint
benchmarks of \cref{sec:jointprior} do not, comparing it directly to their
figures would understate us. We therefore also report
\emph{correct-and-secure}: correctness conjoined with safety, omitting the
equivalence term, which is the construction those benchmarks use.

\begin{table}[htbp]
\centering
\footnotesize
\setlength{\tabcolsep}{4pt}
\input{tab_crossbench}
\caption{Our correct-and-secure rate against joint correctness-security
results from other domains. Task types, vulnerability classes and evaluation
methods all differ substantially. The comparison is one of magnitude and of
the ratio between functional and joint success, not of models or difficulty.}
\label{tab:crossbench}
\end{table}

\textbf{We do not reproduce the pattern.} \citet{vero2025baxbench} report that roughly half of functionally correct
backends are exploitable. \citet{chen2025secureagentbench} report 15.2\%
correct-and-secure for their strongest agent, against substantially higher
functional correctness. \citet{patir2025dualgauge} report secure-pass@1 below
12\% while functional pass@1 exceeds 50\%. In each case the security criterion removes a large
majority of functionally correct output.

In our setting it removes 7.4\%.

We think the divergence is more useful than agreement would have been. Three
explanations seem plausible. We cannot separate them with this data:

\begin{enumerate}
\item \textbf{Attack surface.} A regular expression is a single expression
with one failure mode of interest. A backend application has many
independently exploitable components, so the probability that at least one is
vulnerable compounds with program size.
\item \textbf{Vulnerability class.} ReDoS is a structural property of the
pattern, plausibly well represented in training data as a recognised
anti-pattern. CWE-classified defects in application code are more diverse and
more context-dependent.
\item \textbf{Task difficulty ceiling.} Our functional pass rate is itself low
(38--47\%), so the correct subset may be biased toward simpler tasks whose
solutions have less structural room for catastrophic backtracking.
\end{enumerate}

So the claim that generalises is weaker than ``correct code is often
insecure,'' and more useful. \emph{The size of the correctness-to-security
penalty depends on the domain and has to be measured there.} Assuming it
transfers is what we did, and we were wrong. Worse, our composite showed a gap
of about the expected size for entirely the wrong reason
(\cref{sec:decomp}), which is exactly the situation in which an undecomposed
joint metric will mislead you.

\subsection{Cost}

\begin{table}[htbp]
\centering
\small
\begin{tabular}{lrrr}
\toprule
Model & \usable{}@3 (\%) & cost/task (\$$\times 10^{-6}$) & total (\$) \\
\midrule
\texttt{deepseek-v4-flash-0731} & 19.8 & 25.6 & 0.03 \\
\texttt{gpt-5.6-luna} & 18.7 & 42.9 & 0.06 \\
\texttt{gemini-3.1-flash-lite} & 17.1 & 90.2 & 0.12 \\
\texttt{qwen3.6-plus} & 19.8 & 120.6 & 0.16 \\
\texttt{glm-5.2} & 18.7 & 158.2 & 0.21 \\
\texttt{qwen3.6-max-preview} & 21.6 & 387.9 & 0.52 \\
\texttt{gpt-5.6-terra} & 18.7 & 406.2 & 0.55 \\
\texttt{claude-sonnet-5} & 18.0 & 932.4 & 1.26 \\
\texttt{kimi-k3} & 24.8 & 1328.1 & 1.79 \\
\texttt{gpt-5.6-sol} & 21.1 & 2107.6 & 2.85 \\
\texttt{claude-opus-5} & 23.0 & 2514.2 & 3.39 \\
\bottomrule
\end{tabular}

\caption{Cost per task at $n=3$, measured from per-request billing returned
by the serving layer rather than estimated from list prices.}
\label{tab:cost}
\end{table}

The extremes differ by a factor of 98 in price and about three percentage
points in \usable{}@3, which is at the boundary of what we can resolve
(\cref{sec:paired}). This observation is about this task only. It is not a
general claim about model capability.

\subsection{Refusals}

A content filter blocked \texttt{claude-opus-5} on 29 requests, 2.1\% of its
calls, spread across five tasks. Four have some security flavour: strings that
exclude a single quotation mark, a six-character alphanumeric password, hex
byte sequences, and SQL update and insert statements. The fifth is
\emph{``Matches a file extention.''} No other model refused anything.

The refusals were not deterministic. Several tasks came back refused on one
sample and answered on the other two, under identical settings. You can only
see that at $k>1$. A single-sample protocol would have logged a flat failure
and moved on. We count refusals as their own failure category, not as wrong
answers (\cref{app:failures}).

\subsection{Inference-time compute}

On the corpus's simplest task (\emph{``Matches exactly 1 numeric digit
(0-9).''}), \texttt{qwen3.6-max-preview} emitted 1{,}571 hidden reasoning
tokens before answering \texttt{\^{}[0-9]\$}, at a cost of \$0.0098. With
reasoning disabled it produced the identical answer in 10 tokens at
$1/100$th the cost. Asking that model for \emph{low} reasoning effort produced \emph{2{,}375}
reasoning tokens, more than the default. Effort controls are not reliably
monotone.

A reasoning-enabled comparison on 12 tasks across five models yielded
per-request costs $15.7\times$ higher and no interpretable accuracy signal at
that sample size ($\pm14$ percentage points per cell). We report the cost
result and explicitly decline to report an accuracy conclusion.

\section{Measurement Validity}
\label{sec:validity}

\subsection{The reference standard contains errors}
\label{sec:goldstandard}

First we ruled out a loading fault. 449 of 450 reference patterns satisfy
their own tests, so the corpus semantics and dialect are configured correctly.
The real defect is quieter. A reference can pass its own tests and still fail
to denote what its description says, whenever the tests never exercise the
difference.

We drew a seeded random sample of 60 cases in which a model \emph{passed
every test} yet received a \textsc{diff} verdict, and adjudicated 14 in
detail.

\begin{table}[htbp]
\centering
\small
\begin{tabular}{lrr}
\toprule
Adjudicated cause & count & share \\
\midrule
Reference pattern is incorrect & 5 & 36\% \\
Description underspecifies the disputed property & 6 & 43\% \\
Model output is incorrect & 3 & 21\% \\
\bottomrule
\end{tabular}
\caption{Adjudication of 14 sampled \textsc{diff} verdicts on test-passing
outputs. Per-case reasoning is released.}
\label{tab:adjudication}
\end{table}

Separately, 19 of the 60 sampled disagreements (32\%) differ only on
non-ASCII input. Python's \texttt{\textbackslash d} is Unicode-aware and
\texttt{[0-9]} is not, so the two are formally distinct languages and
practically the same thing. Put both effects together and the model is clearly
at fault in roughly \textbf{15\%} of \textsc{diff} verdicts.

Representative cases:

\begin{itemize}
\item \textbf{Description:} \emph{``It just accepts only positive numbers.''}
The reference \lstinline!^\d+([.,]?\d+)?$! admits \texttt{0}. The model
excluded zero and was scored incorrect.

\item \textbf{Description:} \emph{``A very simple ISBN validation
expression---it just checks for a 10 digit number.''} The reference is
\lstinline!^\d{9}[\d|X]$!. The character class contains digit, \emph{pipe}
and \texttt{X}: an alternation operator written inside a class, where it
denotes a literal. The reference therefore accepts \texttt{000000000|}, which
is the witness separating it from the model's \lstinline!^\d{9}[\dX]$!.

\item \textbf{Description:} \emph{``\ldots make sure commas are in the rite
place (if present).''} The reference
\begin{lstlisting}
^\$?\d{1,3}(,?\d{3})*(\.\d{1,2})?$
\end{lstlisting}
makes the comma optional, admitting \texttt{0,000000} and negating the stated
purpose. The model enforced comma placement and was scored incorrect.

\item \textbf{Underspecification:} \emph{``Matches any single upper- or
lower-case letter.''} The reference \lstinline!^[a-zA-Z]$! requires a
whole-string match. A candidate \lstinline![A-Za-z]! requires only
occurrence. The description does not distinguish the two. We tested whether anchoring alone
explains failures generally: it accounts for approximately 6\%, so the effect
is diffuse rather than a single correctable convention.
\end{itemize}

\paragraph{Consequence.} \dfaeq{} and \usable{} are lower bounds on model
correctness by an unmodelled margin. \pass{} and \vuln{} are unaffected,
since neither consults the reference.

\paragraph{Limitation of this analysis.} Fourteen cases, judged by us, on our
own benchmark. That is a weak basis for a precise number and we do not want it
cited as one. We release the reasoning for each case, and the unadjudicated
cases keep an empty verdict field so someone else can fill them in. Anyone
quoting 15\% should re-derive it from a larger sample with annotators who have
no stake in the answer. We report it because the direction is not in doubt,
and because the direction is what changes how the composite should be read.

\subsection{Paired inference}
\label{sec:paired}

Our first analysis compared systems with independent binomial confidence
intervals. That is the wrong estimator for this design. All eleven models
answer the same items, so task difficulty is a shared random effect. Treating
the systems as independent samples charges every one of them separately for
the same variation.

A paired bootstrap fixes this. Resample tasks with replacement, score every
model on each resample, and form the per-pair differences. Task difficulty
cancels, and only genuine disagreement contributes. With $B = 10{,}000$
resamples and a fixed seed:

\begin{table}[htbp]
\centering
\small
\begin{tabular}{lr}
\toprule
Procedure & pairs resolved at 95\% (of 55) \\
\midrule
Independent binomial intervals & 1 \\
Paired bootstrap, \usable{}@3 & 9 \\
Paired bootstrap, \pass{}@3 & 15 \\
Paired bootstrap, \vuln{}@3 & 15 \\
\bottomrule
\end{tabular}
\caption{Resolvable pairwise comparisons under independent versus paired
inference on identical data.}
\label{tab:paired}
\end{table}

The correction is not marginal. It changes the number of defensible
comparisons by nearly an order of magnitude. We flag it because independent
intervals look like the default in adjacent evaluation work, and the design
they are applied to, all systems answering all items, is close to universal.

Even corrected, the middle of the table does not separate. For example
\texttt{qwen3.6-max-preview} versus \texttt{gpt-5.6-sol} yields
$+0.7\%$ with a 95\% interval of $[-2.9\%, +4.3\%]$.

\subsection{Discriminative power of the corpus}
\label{sec:power}

There is a structural reason behind this. On \usable{}@3, \textbf{62\% of
tasks produce the same outcome for all eleven models}. On \pass{}@3 it is
56\%, on \vuln{}@3, 77\%. Only about 167 of 450 tasks carry any
discriminative signal at all.

For calibration: resolving a 3-percentage-point difference at $p \approx 0.2$
with independent intervals requires roughly 2{,}800 items---more than the
full 762-task corpus contains. Paired inference reduces this requirement
substantially but cannot manufacture signal from items on which every system
agrees. Benchmark construction for system comparison should therefore
optimise for discriminative yield, not corpus size.

\section{Threats to Validity}
\label{sec:threats}

\paragraph{Contamination.} The corpus predates every model we evaluated and
is public. Some of what we measured may be memorisation. We have no held-out
private task set, so we cannot put a bound on it. This is the largest
unaddressed threat in the paper.

\paragraph{Single prompt.} Every result here is conditional on one
unoptimised prompt. \cref{sec:goldstandard} shows how often the task
descriptions leave the anchoring convention unstated, so sensitivity to
phrasing is plausibly large. We did not measure it.

\paragraph{Sampling control.} We do not set temperature
(\cref{sec:metrics}), so it varies by provider default. We checked that
outputs still vary across samples. We did not equalise the sampling.

\paragraph{Estimator degeneracy.} At $n=k=3$ the $\pass{}@k$ estimator
reduces to any-of-3.

\paragraph{Coverage asymmetry.} Nine models cover all 450 tasks.
\texttt{kimi-k3} covers 447, having exhausted our budget mid-run, and
\texttt{claude-opus-5} covers 445 because of refusals. The denominators differ
by under 1\%.

\paragraph{Screening, not proof.} \vuln{} is a lower bound
(\cref{sec:metrics}).

\paragraph{Single corpus.} We do not test whether the observed banding
survives corpus change.

\paragraph{Self-adjudication.} See \cref{sec:goldstandard}.

\section{Recommendations for Evaluation Design}

\begin{enumerate}
\item \textbf{Partition metrics by reference dependence, and decompose any
composite that spans the partition.} Metrics evaluated by execution against
inputs (\pass{}, \vuln{}) are robust to reference-set defects. Metrics
evaluated against a human answer key inherit its error rate. A composite
conjoining both can report a large gap that is almost entirely attributable to
the unreliable term---in our case 87\% of it (\cref{sec:decomp})---while
appearing to be a statement about the reliable ones.

\item \textbf{Sample and adjudicate disagreements.} It costs hours. In our
case it changed the conclusion.

\item \textbf{Use paired inference when all systems answer all items.}
\cref{tab:paired}.

\item \textbf{Report discriminative yield}, not only corpus size
(\cref{sec:power}).

\item \textbf{Pin the serving provider and verify post hoc.} Record what
served each request, not only what was requested (\cref{sec:serving}).

\item \textbf{Instrument the scorer with controls} that must pass and must
fail (\cref{sec:metrics}).

\item \textbf{Release raw generations.} Anyone can then overturn a scoring
decision without buying the inference again. That is the difference between a
result someone can check and one they have to take on faith.
\end{enumerate}

\section{Conclusion}

We set out to measure a gap between correct and shippable regular
expressions. We built a composite metric, it showed a large gap, and
decomposition showed that 87\% of that gap came from the one term we had
already shown to be unreliable. What survives is much smaller. 7.4\% of
functionally correct patterns carry a ReDoS vulnerability, against roughly
50\% of functionally correct backends in comparable work. The
correctness-to-security penalty depends on the domain, and in ours it is
small.

Two findings survive that correction intact, both resting only on criteria
that never consult a human answer key. Models are uniformly safer than the
human-written reference patterns against which they are scored, which locates
ReDoS prevalence in the training distribution rather than in generation. And
the entire eleven-model field lies within eight percentage points across a
$98\times$ range in inference cost.

The methodological result is the one we did not plan on. A metric that
compares against a human answer key inherits that key's error rate, and a
composite built on such a term can produce a finding that is large, plausible
and mostly spurious. We only found this by reading our own output after
collecting it. The audit cost a few hours against a run that cost \$10.95, and
it changed what we were willing to publish. We would do it earlier next
time.

\section*{Acknowledgements}

Scoring is performed by \texttt{regexbench}~\citep{regexbench}, prior work by
the present authors, used here as a pinned dependency and not contributed by
this paper. The corpus is Re(gEx$|$DoS)Eval~\citep{regexeval}, used under its
original terms and not redistributed.

\bibliographystyle{plainnat}
\bibliography{refs}

\appendix

\section{Failure taxonomy}
\label{app:failures}

Of 14{,}850 requests, 54 (0.36\%) failed and are reported rather than
excluded.

\begin{table}[htbp]
\centering
\small
\begin{tabular}{llr}
\toprule
Model & Cause & Count \\
\midrule
\texttt{claude-opus-5} & content-filter refusal & 29 \\
\texttt{kimi-k3} & account spending limit reached mid-run & 11 \\
\texttt{kimi-k3} & response without resolved provider & 3 \\
\texttt{gpt-5.6-sol} & response without resolved provider & 1 \\
\bottomrule
\end{tabular}
\caption{Request failures by cause. Responses lacking a resolved provider are
discarded rather than scored, as provenance is a precondition for
reproducibility.}
\end{table}

Empty completions are classified by reported \texttt{finish\_reason} into
refusal, token-budget exhaustion, and unexplained, rather than being scored
as empty patterns---which would be indistinguishable from a maximally poor
answer.

\section{Operational hazards}
\label{app:hazards}

We document configuration hazards encountered during this study, each of
which produced plausible-looking but invalid data. We believe this list has
independent value.

\begin{enumerate}
\item \textbf{Parameter combinations that are silently unroutable.} Setting
\texttt{temperature} together with \texttt{require\_parameters=true} yields a
hard routing failure on 6 of 11 models. Our first pilot failed all 396 calls.

\item \textbf{Systems that cannot satisfy the experimental condition.} Two
Gemini variants return \emph{``Reasoning is mandatory for this endpoint''}
and cannot be evaluated with reasoning disabled. Undetected, this would have
placed one reasoning-enabled system in an otherwise reasoning-disabled
comparison.

\item \textbf{Empty completions from budget exhaustion.} At a 200-token cap,
reasoning models consumed the entire budget on hidden reasoning and returned
no content.

\item \textbf{Resumption across configuration changes.} Our collection is
resumable by (model, task, sample). An aborted run at one token cap left
truncated rows that a resumed run at a different cap treated as complete,
interleaving two configurations indistinguishably. Rows now carry a
configuration fingerprint and resumption refuses on mismatch.

\item \textbf{Sample collapse in scoring.} An early scorer retained only the
final sample per task, which would have rendered all $@3$ statistics
disguised $@1$ statistics with no visible symptom.
\end{enumerate}

\section{Reproduction}

\begin{lstlisting}[language=bash]
git clone https://github.com/foothills-labs/regexleaderboard
cd regexleaderboard
make setup      # pinned scorer + corpus download
make score RUN=sweep
\end{lstlisting}

Scoring reads only the committed generations. No API access, no
expenditure. We verified reproduction from a clean checkout with fresh
dependency installation and corpus download, obtaining bit-identical metrics.
A reduced form of this check executes in continuous integration on every
commit.

\end{document}
